\section{Space Complexity --- Machines --- Lib}
%%% generated by the blueprint pipeline from TCSlib.Complexity.SpaceComplexity.Machines.Lib
\subsection*{Space Complexity --- Machines --- Lib: Introduction}

Declarations of the Lean module \texttt{TCSlib.Complexity.SpaceComplexity.Machines.Lib}.

\subsection*{Space Complexity --- Machines --- Lib: Results}

\begin{lemma}[Ordinary step at a non-call state]
\lean{Complexity.LogProg.rstep_noncall}
\label{Complexity.LogProg.rstep_noncall}
\leanok
\uses{Complexity.LogProg.RProg, Complexity.LogProg.rstep, Turing.Cfg, Turing.MultiTapeTM.step}
\difficulty{1}
Let $P$ be a register program with oracle, and let $c$ be a configuration whose state is
$l$. If $P$ has no oracle call at $l$, then the register-program step from $c$ equals
the ordinary machine step of $P$ from $c$.
\end{lemma}

\begin{lemma}[One step of a run]
\lean{Complexity.LogProg.rrun_one}
\label{Complexity.LogProg.rrun_one}
\leanok
\uses{Complexity.LogProg.RProg, Complexity.LogProg.rrun, Complexity.LogProg.rstep, Turing.Cfg}
\difficulty{1}
Let $P$ be a register program, let $\mathcal O$ be an oracle assigning a Boolean answer
to each decider and Boolean word, and let $c$ be a configuration. Running the program
for one step from $c$ gives its one-step successor configuration.
\end{lemma}

\begin{lemma}[Zero-step run]
\lean{Complexity.LogProg.rrun_zero}
\label{Complexity.LogProg.rrun_zero}
\leanok
\uses{Complexity.LogProg.RProg, Complexity.LogProg.rrun, Turing.Cfg}
\difficulty{1}
For every register program $P$ with oracle and every configuration $c$, the run of
length zero from $c$ is $c$ itself.
\end{lemma}

\begin{lemma}[Addition of run times]
\lean{Complexity.LogProg.rrun_add}
\label{Complexity.LogProg.rrun_add}
\leanok
\uses{Complexity.LogProg.RProg, Complexity.LogProg.rrun, Turing.Cfg}
\difficulty{3}
Let $P$ be a register program, let $\mathcal O$ be an oracle assigning a Boolean answer
to each decider and Boolean word, and let $c$ be a configuration. For all natural
numbers $a,b$, running the program for $a+b$ steps from $c$ gives the same configuration
as running it for $a$ steps from $c$ and then for $b$ more steps.
\end{lemma}

\begin{definition}[Regular-register configuration]
\lean{Complexity.LogProg.regCfg}
\label{Complexity.LogProg.regCfg}
\leanok
\uses{Turing.Cfg}
Let $c$ be a configuration of a machine with $m$ Boolean work tapes, let $s$ be a state,
let $r$ be one of the work-tape indices, and let
$f\colon\mathbb{Z}\to\mathrm{Option}(\mathrm{Bool})$ and $p\in\mathbb{Z}$. The
regular-register configuration has state $s$, the same input-head position and output as
$c$, tape contents $f$ and head position $p$ on tape $r$, and the same contents and head
positions as $c$ on every other work tape.
\end{definition}

\begin{lemma}[State of a register configuration]
\lean{Complexity.LogProg.regCfg_state}
\label{Complexity.LogProg.regCfg_state}
\leanok
\uses{Complexity.LogProg.regCfg, Turing.Cfg}
\difficulty{1}
Let $c$ be a configuration, let $s$ be a state, let $r$ be a register, let $f$ be an
integer-indexed tape of optional Boolean symbols, and let $p$ be an integer. The
configuration obtained by replacing register $r$'s tape by $f$, placing its head at $p$,
and setting the state to $s$ has state $s$.
\end{lemma}

\begin{definition}[Regular-register action]
\lean{Complexity.LogProg.regAct}
\label{Complexity.LogProg.regAct}
\leanok
\uses{Turing.Action}
Let $r$ be one of $m$ Boolean work tapes, let $w$ be an optional instruction for writing
on that tape, let $v\in\{-1,0,1\}$ be its head movement, and let $s$ be a state. The
corresponding machine action leaves the input head stationary, applies $(w,v)$ on tape
$r$, leaves every other work tape unchanged with its head stationary, emits no output
symbol, and enters state $s$.
\end{definition}

\begin{lemma}[A register action on a configuration]
\lean{Complexity.LogProg.apply_regAct}
\label{Complexity.LogProg.apply_regAct}
\leanok
\uses{Complexity.LogProg.regAct, Complexity.LogProg.regCfg, Turing.Action.apply, Turing.Cfg, Turing.moveInputPos_zero}
\difficulty{4}
Let $c$ be a configuration with $m$ Boolean work tapes, state set $\Lambda$, and fixed
input $x$. Given states $s,s'\in\Lambda$, a register $r\in\{0,\ldots,m-1\}$, a tape
function $f:\mathbb Z\to\mathrm{Option}(\{\mathrm{false},\mathrm{true}\})$, a position
$p\in\mathbb Z$, a write instruction
$w\in\mathrm{Option}(\mathrm{Option}(\{\mathrm{false},\mathrm{true}\}))$, and movement
$mv\in\{-1,0,1\}$, applying the action that performs $w$ on register $r$, moves its head
by $mv$, and enters $s'$ to the configuration obtained from $c$ by setting its state to
$s$, register $r$'s tape to $f$, and its head position to $p$ gives the configuration
with state $s'$, tape $f$ if $w$ is absent and otherwise $f$ updated at $p$ with the
specified optional symbol, and head position $p+mv$. All other configuration data are
unchanged.
\end{lemma}

\begin{lemma}[Reading a register after replacement]
\lean{Complexity.LogProg.regCfg_read}
\label{Complexity.LogProg.regCfg_read}
\leanok
\uses{Complexity.LogProg.regCfg, Turing.Cfg, Turing.Cfg.workTapeSymbols}
\difficulty{1}
Let $c$ be a configuration, let $s$ be a state, let $r$ be a register, let $f$ assign an
optional Boolean symbol to each integer tape cell, and let $p$ be an integer. In the
configuration obtained from $c$ by setting the state to $s$, replacing register $r$'s
tape by $f$, and placing its head at $p$, the symbol scanned by register $r$ is $f(p)$.
\end{lemma}

\begin{lemma}[Updating a Tape at a Word Position]
\lean{Complexity.LogProg.update_bufferTape_cons}
\label{Complexity.LogProg.update_bufferTape_cons}
\leanok
\uses{Turing.FinTM.bufferTape}
\difficulty{4}
Using juxtaposition for concatenation, let $\mathrm{pre}$ and $w$ be Boolean words and $b,b'$ bits. If a tape stores $\mathrm{pre}bw$, overwriting the cell at position $|\mathrm{pre}|$ with $b'$ makes it store $\mathrm{pre}b'w$.
\end{lemma}

\begin{lemma}[Appending a Bit to a Tape]
\lean{Complexity.LogProg.update_bufferTape_end}
\label{Complexity.LogProg.update_bufferTape_end}
\leanok
\uses{Turing.FinTM.bufferTape, Turing.FinTM.bufferTape_append}
\difficulty{1}
For any Boolean word $w$ and bit $b$, overwriting the tape storing $w$ at cell $|w|$
with $b$ gives exactly the tape storing $w$ followed by $b$.
\end{lemma}

\begin{definition}[Increment-counter action]
\lean{Complexity.LogProg.incCAct}
\label{Complexity.LogProg.incCAct}
\leanok
\uses{Complexity.LogProg.regAct, Turing.Action}
Let $r$ be one of $m$ Boolean work tapes, let $c_C,c_B$ be states, and let $rd$ be an
optional Boolean read from tape $r$. If $rd$ is the bit $1$, the action writes $0$ on
tape $r$, moves its head one cell right, and enters $c_C$. In every other case, it
writes $1$ on tape $r$, moves its head one cell left, and enters $c_B$.
\end{definition}

\begin{definition}[Decrement-counter action]
\lean{Complexity.LogProg.incBAct}
\label{Complexity.LogProg.incBAct}
\leanok
\uses{Complexity.LogProg.regAct, Turing.Action}
Let $r$ be one of $m$ Boolean work tapes, let $c_B$ and $n$ be states, and let $rd$ be
an optional Boolean read from tape $r$. If a bit is present, the action leaves tape $r$
unchanged and moves its head one cell left, entering $c_B$. If no bit is present, it
leaves tape $r$ unchanged and moves its head one cell right, entering $n$.
\end{definition}

\begin{lemma}[Run of the increment carry-in phase]
\lean{Complexity.LogProg.incC_run}
\label{Complexity.LogProg.incC_run}
\leanok
\uses{Complexity.LogProg.apply_regAct, Complexity.LogProg.incCAct, Complexity.LogProg.incW, Complexity.LogProg.regCfg, Complexity.LogProg.regCfg_read, Complexity.LogProg.regCfg_state, Complexity.LogProg.rrun, Complexity.LogProg.rrun_add, Complexity.LogProg.rrun_one, Complexity.LogProg.rstep_noncall, Complexity.LogProg.update_bufferTape_cons, Complexity.LogProg.update_bufferTape_end, Turing.Cfg, Turing.FinTM.bufferTape, Turing.MultiTapeTM.step}
\difficulty{6}
Using juxtaposition for concatenation, let $P$ be a register program with an oracle, and suppose its transition at label $c_C$ is the carry-in action for incrementing register $r$, with the next state determined by the scanned bit; assume $c_C$ is not an oracle-call label. For any configuration $c$ and Boolean words $w$ and $\mathrm{pre}$, there are a time $T$ and head position $p$ with $-1\le p<|\mathrm{pre}\,\mathrm{incW}(w)|$ such that the run from state $c_C$, with tape contents $\mathrm{pre}w$ and head at $|\mathrm{pre}|$, reaches state $c_B$ with tape contents $\mathrm{pre}\mathrm{incW}(w)$ and head at $p$. At every earlier time, the state is $c_C$, the tape contains some word $f$, and the head position $q$ satisfies $|\mathrm{pre}|\le q\le|\mathrm{pre}w|$.
\end{lemma}

\begin{lemma}[Run of the increment carry phase]
\lean{Complexity.LogProg.incB_run}
\label{Complexity.LogProg.incB_run}
\leanok
\uses{Complexity.LogProg.apply_regAct, Complexity.LogProg.incBAct, Complexity.LogProg.regCfg, Complexity.LogProg.regCfg_read, Complexity.LogProg.regCfg_state, Complexity.LogProg.rrun, Complexity.LogProg.rrun_add, Complexity.LogProg.rrun_one, Complexity.LogProg.rstep_noncall, Turing.Cfg, Turing.FinTM.bufferTape, Turing.MultiTapeTM.step}
\difficulty{6}
Let $P$ be a register program with oracle, and suppose its transition at label $c_B$ is
the increment-carry action on register $r$, continuing at $c_B$ when the scanned symbol
is present and otherwise proceeding to the next label. Suppose also that $c_B$ is not an
oracle-call label. For any configuration $c$, word $w$, natural number $n$, and integer
$p$ with $p+1=n$ and $p<|w|$, starting with the head at $p$ on a tape containing $w$
reaches after $n+1$ steps the same tape with the head at $0$ and state `next`. At every
earlier time $t<n+1$, the configuration has state $c_B$, the same tape, and some head
position $q$ with $-1\le q\le p$.
\end{lemma}

\begin{lemma}[Increment run on a binary counter]
\lean{Complexity.LogProg.inc_run}
\label{Complexity.LogProg.inc_run}
\leanok
\uses{Complexity.LogProg.bits_succ, Complexity.LogProg.incB_run, Complexity.LogProg.incC_run, Complexity.LogProg.incW, Complexity.LogProg.regCfg, Complexity.LogProg.rrun, Complexity.LogProg.rrun_add, Turing.Cfg, Turing.FinTM.bufferTape}
\difficulty{5}
Let $P$ be a register program with oracle, let $r$ be a register, and let
$c_C,c_B,	ext{next}$ be states. Suppose the transition at $c_C$ is always the action
denoted by `incCAct` for register $r$, and the transition at $c_B$ is always the action
denoted by `incBAct` for register $r$; suppose neither state is a call state. From any
configuration $c$, there is a finite run from the configuration with state $c_C$,
register $r$ holding the binary word of $n$, and head position zero to the configuration
with state $	ext{next}$, register $r$ holding the binary word of $n+1$, and head
position zero. At every earlier time $t$, the configuration has the form of $c$ with
register $r$ holding some tape word and head position $q$, its state is either $c_C$ or
$c_B$, and $-1\le q\le |\mathrm{bits}(n+1)|$.
\end{lemma}
